% ---------------------------------------------------------------
\section{Guarantees and Correctness Properties}
\label{sec:guarantees}
% ---------------------------------------------------------------

The Bailout Protocol's correctness is expressed as three fundamental
guarantees---SAFETY, LIVENESS, and ACCOUNTABILITY---derived from the state
machine $\mathcal{B}$'s structural invariants and the \fact{} layer's pre-commit
enforcement semantics.  These are not operational aspirations; they are
architectural properties that hold whenever: (i) the \fact{} layer enforces
transitions as specified, (ii) the parent-pointer chain $\alpha(\cdot)$ is
computationally intact, and (iii) the human anchor $h(n)$ is provisioned at
node initialization.  Each guarantee is stated as a theorem with a proof sketch
tracing the relevant transition structure.

% ---------------------------------------------------------------
\subsection{Guarantee 1: SAFETY --- No Autonomous Collapse}
\label{sec:guarantee-safety}

\begin{definition}[Safety Property (SAFETY)]
\label{def:safety}
No node $n$ can remain in a state permitting consequential action without a
human principal on the causal authorization chain.  Equivalently: the system
cannot undergo \emph{autonomous collapse}---a condition in which an agent
continues consequential action indefinitely without a human principal's
knowledge or approval.
\end{definition}

\begin{maintheorem}[SAFETY]
\label{th:safety}
For all nodes $n$ in the spawn tree, for all reachable states $q$ of the Bailout
Protocol state machine $\mathcal{B}$, if $q$ permits consequential action then a
human principal $h \in H$ lies on the causal authorization chain for $n$.
\end{maintheorem}

\begin{Proof}[Proof Sketch]
The proof proceeds by structural analysis of the transition graph of
$\mathcal{B}$, showing that every path to a consequential-action state passes
through a human-credentialed resolution.

\textbf{Base.}  The initial state $q_0 = \texttt{IDLE}$ is reachable only from
\texttt{HUMAN\_ACKNOWLEDGED} via the resolution event
$\sigma_{\text{resolve}}(\texttt{ACK})$ issued by the human anchor $h(n)$
(Transition T9).  This requires authentication against the exclusive human
credential $C_{h(n)}$; no machine actor possesses $C_{h(n)}$.  Therefore,
$q_0$ implies human presence on the authorization chain by construction.

\textbf{Inductive step.}  We examine each non-terminal state and show that no
path leads to a consequential-action state without human resolution.

\begin{itemize}\itemsep=0pt
\item \texttt{BOUNDED}: The pre-commit gate blocks all consequential action.  The
only outgoing transitions are T2 (return to \texttt{IDLE} via human-issued
\texttt{ACK}) and T3 (timeout-forced advance to \texttt{BAIL\_PENDING}).  No
transition from \texttt{BOUNDED} permits consequential action without human
resolution.

\item \texttt{BAIL\_PENDING}: Consequential action is blocked.  The sole
transition T4 is forced by $\tau_{\text{prop}}$ expiration; no machine actor can
suppress or redirect it (SP-2, Section~\ref{sec:bp-state-machine}).

\item \texttt{ESCALATING}: Consequential action is blocked.  The machine exits
via T6 (reaches anchor, transitions to \texttt{HUMAN\_ACKNOWLEDGED}) or T7
(pathway exhausted, terminal failure).  Neither exit permits consequential
action without human acknowledgment.

\item \texttt{HUMAN\_ACKNOWLEDGED}: Consequential action is blocked.  The
machine exits via T8 (human issues binding directive) or T10 (human timeout,
terminal failure).  Only T8 produces a non-terminal accepting state.
\end{itemize}

\textbf{Conclusion.}  The only non-terminal accepting state is \texttt{RESOLVED},
reachable exclusively via T8 with $sigma_{\text{resolve}}$ authenticated against
$C_{h(n)}$.  \texttt{IDLE} (the state permitting consequential action) is
reachable only by returning through this same human-credentialed path.  No cycle
exists in the transition graph that permits a machine actor to re-authorize
consequential action after the state machine has exited \texttt{IDLE} without
human intervention.  $\square$
\end{Proof}

The SAFETY guarantee is structurally enforced by SP-1 (non-regression) and
SP-3 (human-only resolution) from the state machine definition.  It does not
depend on any machine actor choosing to comply; it follows from the deterministic
transition relation alone.

% ---------------------------------------------------------------
\subsection{Guarantee 2: LIVENESS --- Eventual Human Contact}
\label{sec:guarantee-liveness}

\begin{definition}[Liveness Property (LIVENESS)]
\label{def:liveness}
Every exception state $e$ that triggers the Bailout Protocol at a node $n$
reaches the designated human anchor $h(n)$ in finite time.  Formally: for all
nodes $n$ and all exception states $e$ satisfying the trigger condition
$\text{TC}(n,x)$, the escalation path $C(n)$ generated by
Algorithm~\ref{alg:escalation} terminates at $h(n)$ within a bounded number of
state machine transitions.
\end{definition}

\begin{maintheorem}[LIVENESS]
\label{th:liveness}
For all nodes $n$ and all exception states $e$ at $n$, the Bailout Protocol
state machine reaches the state \texttt{HUMAN\_ACKNOWLEDGED} within at most
$T_{\max}$ transitions, where
$T_{\max} = 4 + \bigl(d(n) - d(h(n))\bigr)$ and $d(\cdot)$ denotes depth in the
spawn tree.
\end{maintheorem}

\begin{Proof}[Proof Sketch]
The proof enumerates the transition sequence from trigger condition firing to
human acknowledgment and shows that each segment contributes a bounded number of
transitions.

\begin{itemize}\itemsep=0pt
\item \textbf{Trigger to escalation (bounds timeout).}  When TC fires at $n$,
the state machine transitions IDLE $\rightarrow$ BOUNDED (T1) and
$\tau_{\text{bounds}}$ is armed.  By SP-2 (timeout-forced progression), the
machine transitions to \texttt{BAIL\_PENDING} (T3) when $\tau_{\text{bounds}}$
expires, regardless of any machine actor's assessment.  This first forced
transition occurs within at most $T_{\text{bounds}}$ time units.

\item \textbf{Escalation initiation (propagation timeout).}  From
\texttt{BAIL\_PENDING}, the sole transition is to \texttt{ESCALATING} (T4),
forced by $\tau_{\text{prop}}$ expiration.  No machine actor can suppress this
transition.  The state enters \texttt{ESCALATING} within at most
$T_{\text{prop}}$ additional time units.

\item \textbf{Propagation along the parent chain.}  The escalation loop
(Algorithm~\ref{alg:escalation}, T5 self-loop) advances the propagation path one
hop per successful \texttt{SendToParent} invocation.  Each hop moves from a node
$v$ to its parent $\alpha(v)$, strictly decreasing depth:
$d(\alpha(v)) = d(v) - 1$ for all $v$ with $v \neq h(n)$.  The chain terminates
at $h(n)$ where depth is minimal.  The maximum number of hops is
$d(n) - d(h(n)) \leq D_{\max}$, a finite structural bound established at root
initialization.  Each hop is bounded by $\tau_{\text{prop}}$; retries do not
extend the total hop count beyond $D_{\max}$.

\item \textbf{Human acknowledgment.}  When propagation reaches the anchor node,
the state machine transitions to \texttt{HUMAN\_ACKNOWLEDGED} (T6), completing the
escalation path.
\end{itemize}

\textbf{Combined bound.}  The total number of transitions from the first
trigger to \texttt{HUMAN\_ACKNOWLEDGED} is at most:
\begin{equation}
T_{\max} \;\leq\;
\underbrace{1}_{\text{T1}} \;+\;
\underbrace{1}_{\text{T3}} \;+\;
\underbrace{1}_{\text{T4}} \;+\;
\underbrace{d(n) - d(h(n))}_{\text{T5 loops}} \;+\;
\underbrace{1}_{\text{T6}}
\;=\; 4 + \bigl(d(n) - d(h(n))\bigr).
\end{equation}
Since $d(n) - d(h(n))$ is bounded by $D_{\max}$, a finite constant,
$T_{\max}$ is finite for all nodes.  $\square$
\end{Proof}

The LIVENESS guarantee requires two structural preconditions: the parent-pointer
chain $\alpha(\cdot)$ must be computationally intact, and the spawn tree depth
must be bounded by $D_{\max}$.  Both are established at node initialization and
enforced by the \fact{} layer's pre-commit hook.  Propagation is structurally
compelled at each hop; the guarantee does not require voluntary machine-actor
cooperation.

% ---------------------------------------------------------------
\subsection{Guarantee 3: ACCOUNTABILITY --- Human Always Identifiable}
\label{sec:guarantee-accountability}

\begin{definition}[Accountability Property (ACCOUNTABILITY)]
\label{def:accountability}
Every exception state that triggers the Bailout Protocol produces a Forensic
Ledger entry that (i) uniquely identifies the human principal $h(n)$ responsible
for the exception's resolution, (ii) records the complete causal chain from the
triggering event to the human directive, and (iii) is cryptographically sealed
against post-hoc modification.  The human principal who authorized the
resolution is attributable in the framework's immutable record for all time.
\end{definition}

\begin{maintheorem}[ACCOUNTABILITY]
\label{th:accountability}
For every exception state $e$ that triggers the Bailout Protocol at node $n$,
the Forensic Ledger entry chain generated by the protocol satisfies:
\begin{enumerate}
\item[(i)] the human anchor $h(n)$ is uniquely determined by the spawn tree
structure and is recorded in the \texttt{bailout-trigger} entry at the
originating node;
\item[(ii)] every state transition in the protocol's execution is recorded in
the Forensic Ledger before the transition commits (Invariant I4), forming a
complete causally ordered chain; and
\item[(iii)] the entry chain is SHA-256 hash-chained, making any post-hoc
modification of any entry detectable by hash verification.
\end{enumerate}
\end{maintheorem}

\begin{Proof}[Proof Sketch]
We prove each clause separately.

\medskip
\noindent\textbf{Proof of (i): Unique human anchor identification.}
The human anchor $h(n)$ is provisioned at node $n$'s initialization and stored
in \fact{} layer-managed memory.  It is immutable: the \fact{} layer's
pre-commit hook rejects any runtime attempt to modify $h(n)$ (Invariant I3).
The \texttt{bailout-trigger} Forensic Ledger entry records $h(n)$ explicitly as
part of the structured payload, along with the originating node $n$ and the
trigger condition sub-components.  Because $h(n)$ is established once and never
modified, it is unique and permanent for the node's lifetime.  No machine actor
can substitute a different principal: no machine actor at any plane possesses
the credential provisioning authority required to designate or change $h(n)$.
$\square$

\medskip
\noindent\textbf{Proof of (ii): Complete causal chain.}
Invariant I4 enforces that every state transition in $\mathcal{B}$ produces at
least one corresponding Forensic Ledger entry \emph{before} the transition
commits.  This is enforced by the \fact{} layer's pre-commit hook: any state
transition attempted without a preceding ledger entry is rejected.  The entry
types generated by the protocol (\texttt{bailout-trigger},
\texttt{timeout-bounds-expired}, \texttt{escalation-sent}, \texttt{human-ack},
\texttt{human-directive}) together form a causally ordered chain: each entry
references the node and state at the time of writing, and propagation entries
reference the previous hop's pathway identifier.  The causal ordering is
guaranteed by the pre-commit enforcement timing: the entry is written and
hashed before the state transition commits, so the record of state $q$ is
causally prior to state $q'$.  $\square$

\medskip
\noindent\textbf{Proof of (iii): Cryptographic sealing.}
Each Forensic Ledger entry $E_i$ contains $\hash{E_{i-1}}$, the SHA-256 hash of
the preceding entry $E_{i-1}$, forming the chain
$E_0, E_1, E_2, \ldots, E_k$.  The genesis entry $E_0$ contains its own hash as
a self-sealing bootstrap.  A post-hoc modification of any entry $E_j$ changes
$\hash{E_j}$, which breaks the hash reference in $E_{j+1}$ and all subsequent
entries.  Verification is performed by recomputing the chain hash forward from
$E_0$ and comparing against the stored hash in $E_k$; any discrepancy indicates
tampering.  The hash function is pre-image resistant and collision resistant
under standard cryptographic assumptions, making it computationally infeasible
for a machine actor to modify an entry without detection.  $\square$
\end{Proof}

ACCOUNTABILITY is the conjunction of SAFETY and LIVENESS plus the \fact{}
layer's integrity mechanisms.  SAFETY ensures the human is present on the
authorization chain; LIVENESS ensures the human is contacted; ACCOUNTABILITY
ensures the contact is recorded, attributable, and tamper-evident.  Together they
establish that the \bxthree{} Framework's upstream accountability guarantee is a
formally evidenced property of the runtime execution trace, not a design
aspiration.

% ---------------------------------------------------------------
\subsection{Combined Correctness Statement}
\label{sec:combined-correctness}

\begin{maintheorem}[Combined Correctness]
\label{th:combined-correctness}
For all nodes $n$ in the spawn tree and all exception states $e$ at $n$:
\begin{enumerate}
\item[SAFETY:] the system cannot execute consequential action without a human
principal on the authorization chain;
\item[LIVENESS:] the exception reaches $h(n)$ in finite time; and
\item[ACCOUNTABILITY:] the complete propagation path and the human principal's
resulting directive are recorded in an immutable, cryptographically sealed
Forensic Ledger entry chain.
\end{enumerate}
\end{maintheorem}

The combined statement follows directly from Theorems~\ref{th:safety},
\ref{th:liveness}, and~\ref{th:accountability}.

\begin{conditions-list}[Model Assumptions]
\item The \fact{} layer's pre-commit enforcement operates correctly: state
machine transitions are evaluated and enforced as specified.
\item The parent-pointer chain $\alpha(\cdot)$ is computationally intact: each
node resolves to a valid parent up to the human anchor.
\item The human anchor $h(n)$ is provisioned and reachable at node initialization.
\item The SHA-256 hash function remains pre-image and collision resistant.
\end{conditions-list}

Violations of these assumptions represent failures of the enforcement
substrate, not failures of the protocol itself.  The Forensic Ledger records
evidence of any assumption violation, enabling detection and external auditing.
The formal verification, adaptive timeout, and distributed anchor pool research
agendas (Section~\ref{sec:limitations}) address the conditions under which these
assumptions can be relaxed or strengthened.
